\documentclass[aps,prd,reprint,superscriptaddress,nofootinbib]{revtex4-2}

% ============================================================
% Packages
% ============================================================
\usepackage{amsmath,amssymb,amsfonts,mathtools,bm}
\usepackage{booktabs}
\usepackage{array}
\usepackage{hyperref}
\usepackage{microtype}

\hypersetup{
	colorlinks=true,
	linkcolor=blue,
	citecolor=blue,
	urlcolor=blue
}

% ============================================================
% Useful commands
% ============================================================
\newcommand{\SME}{\mathrm{SME}}
\newcommand{\DM}{\mathrm{DM}}
\newcommand{\FW}{\mathrm{FW}}
\newcommand{\CPT}{\mathrm{CPT}}
\newcommand{\ACF}{\mathrm{ACF}}
\newcommand{\NR}{\mathrm{NR}}
\newcommand{\mbf}{\mathbf}
\newcommand{\msf}{\mathsf}

% ============================================================
% Document
% ============================================================
\begin{document}
	
	\title{%
		Relativistic--Nonrelativistic Matching and Matter--Antimatter
		Constraints for Exotic Spin-Dependent Interactions:
		A Unified Framework and the Limits of One-Sided Bounds as CPT Tests
	}
	
	\author{Temidayo Solomon Oyewo}
	\author{Oluwole Emmanuel Oyewande}
	\email{oe.oyewande@ui.edu.ng}
	\affiliation{Department of Physics, University of Ibadan, Ibadan, Nigeria}
	
	\date{\today}
	
	% ============================================================
	% Abstract
	% ============================================================
	\begin{abstract}
		
		Light, weakly coupled bosons can mediate exotic spin-dependent
		interactions over a wide range of distances. Experimental constraints on
		such interactions are commonly reported using the nonrelativistic
		potential basis originating with Dobrescu and Mocioiu (DM), whereas
		Lorentz- and CPT-violating effects are naturally formulated within the
		Standard-Model Extension (SME). We establish a controlled correspondence
		between selected coefficients in the minimal SME fermion sector and
		nonrelativistic operator structures by applying the field normalisation
		and Foldy--Wouthuysen (FW) procedure appropriate to derivative
		kinetic terms. The analysis is restricted to $b_\mu$, $H_{\mu\nu}$, and
		$d_{\mu\nu}$ and does not claim a complete SME--DM dictionary. The
		canonical nonrelativistic Hamiltonian reproduces the established
		spin and spin--momentum structures, including the mass-enhanced
		$d_{i0}$ contribution and the momentum-linear combination
		$d_{ij}+d_{00}\delta_{ij}$.
		
		We then investigate whether published matter and antimatter upper limits
		can be converted into a quantitative CPT asymmetry. A signed quantity
		such as
		$A=(g_f-g_{\bar f})/(g_f+g_{\bar f})$
		is meaningful only when $g_f$ and $g_{\bar f}$ are signed observables or
		well-defined estimators. Experimental exclusion results are instead
		usually one-sided bounds on positive quantities such as $|g|$. If
		$U_f=s_f g_\star$ and $U_{\bar f}=s_{\bar f}g_\star$ describe two such
		limits on the same underlying coupling magnitude, their apparent
		asymmetry is exactly
		\[
		A_U=\frac{U_f-U_{\bar f}}{U_f+U_{\bar f}}
		=\frac{s_f-s_{\bar f}}{s_f+s_{\bar f}},
		\]
		which contains no CPT parameter. We therefore interpret this quantity
		as a sensitivity-gap diagnostic rather than as a CPT observable.
		
		A compiled registry contains 273 constraint records, of which 258 are
		admitted to the present analysis. Nineteen analysed records involve
		antimatter, and eleven matter--antimatter pairs satisfy the adopted
		matching requirements in coupling convention, operator label, fermion
		sector, and interaction range. The resulting bound-based asymmetries
		span approximately $0.33$--$1$ in magnitude, but their size is not
		evidence for CPT violation. The present coverage is strongly
		matter-dominated, with substantial deficits in electron--nucleus,
		neutron--nucleus, and proton--nucleus sectors. We formulate a practical
		two-stage programme in which balanced sensitivity regions are first
		identified through a sensitivity-ratio function and genuine CPT
		tests are subsequently performed using signed observables or full
		likelihood functions.
		
	\end{abstract}
	
	\maketitle
	
	% ============================================================
	\section{Introduction}
	% ============================================================
	
	Searches for weakly coupled new interactions at low energies provide a
	complementary route to high-energy searches for physics beyond the
	Standard Model (BSM). Light bosons, including axionlike particles and
	new spin-1 states, can mediate long-range forces whose effects may be
	probed through precision spectroscopy, atomic and molecular systems,
	torsion balances, comagnetometers, spin sensors, and antimatter
	experiments \cite{MoodyWilczek1984,DobrescuMocioiu2006,
		Cong2025,WuYan2026}.
	
	A standard phenomenological language for such interactions is provided
	by the sixteen rotationally invariant nonrelativistic structures
	introduced by Dobrescu and Mocioiu \cite{DobrescuMocioiu2006}.
	Modern treatments have clarified the relation between different
	historical potential conventions, reduced coupling constants, and
	experimental applications \cite{Cong2025}. The DM classification is,
	however, a nonrelativistic phenomenological framework rather than a
	complete relativistic effective field theory. In particular, the label
	of a DM potential does not by itself determine the Lorentz or CPT
	properties of an underlying interaction.
	
	The SME provides a systematic effective-field-theory description of
	Lorentz and CPT violation \cite{ColladayKostelecky1997,
		ColladayKostelecky1998}. In its minimal fermion sector, coefficients
	such as $b_\mu$, $H_{\mu\nu}$, and $d_{\mu\nu}$ modify the Dirac
	Lagrangian and produce observable nonrelativistic effects. The
	Foldy--Wouthuysen transformation provides the standard route from the
	relativistic Dirac description to a two-component nonrelativistic
	Hamiltonian \cite{FoldyWouthuysen1950,KosteleckyLane1999}.
	
	The relationship between these descriptions is useful but must be
	formulated carefully. An SME coefficient is a one-fermion background
	quantity, whereas a DM coupling such as $g_A^e g_A^N$ is a product of
	couplings at two interaction vertices. Consequently, an SME coefficient
	cannot in general be identified numerically with a DM coupling.
	A legitimate connection must proceed through the nonrelativistic
	operator structure and, where required, through an explicit
	two-fermion scattering amplitude.
	
	This distinction is especially important for matter--antimatter
	comparisons. Antiprotonic helium spectroscopy has already constrained
	exotic spin-dependent interactions involving antiprotons
	\cite{Ficek2018}, while positronium provides a theoretically clean
	leptonic matter--antimatter system for precision tests
	\cite{Adkins2022}. At the same time, spin-based quantum sensors and
	NV-diamond techniques continue to improve constraints on exotic
	interactions in ordinary matter \cite{Jiao2021,Liang2023,Jiang2024}.
	These developments make a systematic comparison of matter and
	antimatter constraints timely.
	
	A natural dimensionless quantity for signed couplings is
	\begin{equation}
		A_\alpha =
		\frac{g_\alpha^f-g_\alpha^{\bar f}}
		{g_\alpha^f+g_\alpha^{\bar f}},
		\label{eq:Adef}
	\end{equation}
	where $f$ and $\bar f$ denote the corresponding matter and antimatter
	sectors. Equation~\eqref{eq:Adef} is meaningful when the two quantities
	are signed measurements or estimators with well-defined probability
	distributions. Most searches for weak exotic interactions instead
	publish one-sided exclusion limits on positive quantities such as
	$|g_\alpha|$. Applying Eq.~\eqref{eq:Adef} directly to such limits
	therefore changes its interpretation.
	
	The central statistical result of this work is that the resulting
	quantity measures relative experimental sensitivity rather than
	CPT violation. This follows algebraically and does not depend on the
	details of the experimental data set.
	
	The work has three objectives. First, we establish a controlled
	relativistic--nonrelativistic correspondence for selected SME fermion
	coefficients. Second, we distinguish genuine signed matter--antimatter
	tests from comparisons of one-sided exclusion limits. Third, we use a
	compiled constraint registry to quantify present matter--antimatter
	coverage and identify experimental sectors in which matched
	measurements would be particularly informative.
	
	We deliberately do not claim a complete mapping of the SME fermion
	sector onto all sixteen DM potentials. Such a mapping requires additional
	coefficients and, for a true two-body correspondence, explicit matching
	of relativistic amplitudes. Similarly, the numerical asymmetries
	obtained from one-sided upper limits are not interpreted as evidence for
	new physics.
	
	% ============================================================
	\section{Relativistic and nonrelativistic descriptions}
	% ============================================================
	
	\subsection{The SME fermion sector}
	
	Unless otherwise stated, the field-theory equations in this section use
	natural units $\hbar=c=1$ and metric signature $(+,-,-,-)$. To first
	order in the Lorentz-violating coefficients relevant here, the fermion
	Lagrangian may be written as
	\begin{equation}
		\mathcal{L}
		=
		\bar\psi
		\left(
		i\gamma^\mu\partial_\mu-m
		-b_\mu\gamma_5\gamma^\mu
		-\frac{1}{2}H_{\mu\nu}\sigma^{\mu\nu}
		+\frac{i}{2}d_{\mu\nu}
		\gamma_5\gamma^\mu
		\overleftrightarrow{\partial^\nu}
		\right)\psi ,
		\label{eq:SMEL}
	\end{equation}
	where
	\begin{equation}
		\bar\psi\overleftrightarrow{\partial^\nu}\psi
		=
		\bar\psi\,\partial^\nu\psi
		-
		(\partial^\nu\bar\psi)\psi
	\end{equation}
	and
	\begin{equation}
		\sigma^{\mu\nu}
		=
		\frac{i}{2}[\gamma^\mu,\gamma^\nu].
	\end{equation}
	
	The coefficient $b_\mu$ is CPT odd, whereas $H_{\mu\nu}$ and
	$d_{\mu\nu}$ are CPT even in the usual SME classification
	\cite{ColladayKostelecky1997,ColladayKostelecky1998,
		KosteleckyRussell2026}. This classification should not, however, be
	used by itself to infer the sign of a particle--antiparticle energy
	shift. The relation between the particle and antiparticle Hamiltonians
	must be obtained from the complete relativistic theory.
	
	The coefficient $d_{\mu\nu}$ is particularly delicate because it occurs
	in the kinetic sector. Its time-derivative terms must first be removed
	from the canonical time evolution by an appropriate field
	redefinition. Performing a FW reduction before this canonical
	normalisation can lead to spurious sign differences or to a
	non-Hermitian intermediate Hamiltonian \cite{KosteleckyLane1999}.
	
	\subsection{The Dobrescu--Mocioiu framework}
	
	For two nonrelativistic fermions separated by $\mathbf r$, the original
	DM construction identifies sixteen rotationally invariant structures
	formed from spin, relative momentum or velocity, $\hat{\mathbf r}$,
	and the interaction range. For a boson of mass $m_\phi$ the range is
	\begin{equation}
		\lambda=\frac{\hbar}{m_\phi c}.
		\label{eq:lambda}
	\end{equation}
	
	The detailed normalisation of the potentials has appeared in several
	forms in the literature. In particular, the modern review of
	Cong \textit{et al.} gives a complete field-theoretic formulation in
	terms of reduced coupling constants and explicitly discusses the
	differences between historical conventions
	\cite{Cong2025}. We therefore use labels such as $V_2$, $V_{2+3}$,
	$V_{4+5}$, $V_{9+10}$, and $V_{12+13}$ only as database and literature
	labels. We do not reproduce a second set of explicit potential
	formulae here, since doing so would introduce unnecessary convention
	dependence into the SME matching.
	
	This point is important for the present work. The database may contain
	an experimental limit reported for a combined potential such as
	$V_{2+3}$, but that does not justify assigning the result to $V_2$ or
	$V_3$ separately. Combined experimental quantities are therefore
	retained in their published form.
	
	\subsection{One-body SME operators versus two-body potentials}
	
	The correct conceptual relation is
	\begin{equation}
		\text{relativistic SME operator}
		\longrightarrow
		\text{one-body nonrelativistic operator}
		\longrightarrow
		\text{two-body interaction structure},
		\label{eq:matchingchain}
	\end{equation}
	where the final step requires an explicit model for the source or
	mediator.
	
	For example, $b_i$ produces a one-body spin coupling. It is not itself
	the product of two mediator couplings. A two-body interaction involving
	two spins may have a related spin structure, but the numerical
	coefficient depends on the underlying relativistic interaction and on
	the source. The present work therefore establishes structural
	correspondence, not equality of SME coefficients and DM coupling
	constants.
	
	% ============================================================
	\section{Canonical Foldy--Wouthuysen reduction}
	% ============================================================
	
	\subsection{General procedure}
	
	The standard FW transformation separates the positive- and
	negative-energy sectors of the Dirac Hamiltonian and produces a
	two-component Hamiltonian as an expansion in $|\mathbf p|/m$
	\cite{FoldyWouthuysen1950}. For Lorentz-violating kinetic terms, the
	canonical field redefinition must precede the FW transformation
	\cite{KosteleckyLane1999}.
	
	Schematically, after canonical normalisation one writes
	\begin{equation}
		H=\beta m+\mathcal E+\mathcal O,
		\qquad
		[\beta,\mathcal E]=0,
		\qquad
		\{\beta,\mathcal O\}=0,
	\end{equation}
	and the ordinary FW expansion begins with
	\begin{equation}
		H_{\FW}
		=
		\beta m+\mathcal E
		+\frac{\beta\mathcal O^2}{2m}
		-\frac{1}{8m^2}
		[\mathcal O,[\mathcal O,\mathcal E]]
		+\cdots .
		\label{eq:FW}
	\end{equation}
	
	For the free Dirac Hamiltonian,
	\begin{equation}
		\mathcal O_0=\boldsymbol{\alpha}\cdot\mathbf p,
		\qquad
		\mathcal O_0^2=\mathbf p^2,
	\end{equation}
	so that the positive-energy sector begins as
	\begin{equation}
		H_0=m+\frac{\mathbf p^2}{2m}
		+\mathcal O\!\left(\frac{p^4}{m^3}\right).
	\end{equation}
	
	The derivative coefficient $d_{\mu\nu}$ cannot be treated by simply
	replacing $i\partial_\mu$ with a classical momentum in the original
	Lagrangian. The time-derivative part changes the canonical
	normalisation of the field. The controlled treatment is therefore the
	one developed by Kostelecký and Lane \cite{KosteleckyLane1999}.
	
	\subsection{$b_\mu$}
	
	The $b_\mu$ interaction is
	\begin{equation}
		\mathcal L_b
		=
		-b_\mu\bar\psi\gamma_5\gamma^\mu\psi .
	\end{equation}
	
	Retaining the leading nonrelativistic terms from the canonical FW
	Hamiltonian gives
	\begin{equation}
		H_{\NR}^{(b)}
		=
		-b_i\sigma^i
		+\frac{b_0}{m}\,
		\boldsymbol{\sigma}\cdot\mathbf p
		+\mathcal O\!\left(\frac{p^2}{m^2}\right).
		\label{eq:bNR}
	\end{equation}
	
	Thus the spatial component gives a leading spin coupling of order
	$m^0$, whereas $b_0$ contributes to a momentum-linear spin structure
	suppressed by one power of $m$.
	
	The signs in Eq.~\eqref{eq:bNR} are those of the canonical
	nonrelativistic Hamiltonian in the convention used here and agree with
	the general result of Ref.~\cite{KosteleckyLane1999}.
	
	\subsection{$H_{\mu\nu}$}
	
	The tensor interaction is
	\begin{equation}
		\mathcal L_H
		=
		-\frac12 H_{\mu\nu}\bar\psi\sigma^{\mu\nu}\psi .
	\end{equation}
	
	Define
	\begin{equation}
		B_H^i
		=
		\frac12\epsilon^{ijk}H_{jk},
		\qquad
		(H_E)^i=H_{0i}.
	\end{equation}
	
	The leading nonrelativistic Hamiltonian is
	\begin{equation}
		H_{\NR}^{(H)}
		=
		\mathbf B_H\cdot\boldsymbol{\sigma}
		-\frac{1}{m}
		\boldsymbol{\sigma}\cdot
		(\mathbf p\times\mathbf H_E)
		+\mathcal O\!\left(\frac{p^2}{m^2}\right).
		\label{eq:HNR}
	\end{equation}
	
	The first term is a leading spin coupling, while the second is
	linear in momentum. Equation~\eqref{eq:HNR} follows directly from the
	canonical SME Hamiltonian and should not be replaced by a sign inferred
	only from the CPT label of $H_{\mu\nu}$.
	
	\subsection{$d_{\mu\nu}$}
	
	The derivative interaction is
	\begin{equation}
		\mathcal L_d
		=
		\frac{i}{2}
		d_{\mu\nu}
		\bar\psi\gamma_5\gamma^\mu
		\overleftrightarrow{\partial^\nu}\psi .
		\label{eq:Ld}
	\end{equation}
	
	After the required canonical field redefinition and FW transformation,
	the terms relevant through first order in momentum are
	\begin{equation}
		H_{\NR}^{(d)}
		=
		m d_{i0}\sigma^i
		-
		\left(
		d_{ij}+d_{00}\delta_{ij}
		\right)p_j\sigma^i
		+\mathcal O\!\left(\frac{p^2}{m}\right).
		\label{eq:dNR}
	\end{equation}
	
	The mass-enhanced term
	\begin{equation}
		H_{\NR}^{(d_{i0})}
		=
		m d_{i0}\sigma^i
		\label{eq:di0}
	\end{equation}
	is the established result quoted in the SME nonrelativistic
	Hamiltonian \cite{KosteleckyLane1999}.
	
	The momentum-linear contribution is important because the physically
	relevant combination is not $d_{ij}$ or $d_{00}$ separately at this
	order, but
	\begin{equation}
		d_{ij}+d_{00}\delta_{ij}.
		\label{eq:dcombination}
	\end{equation}
	Accordingly, the earlier treatment in which the $d_{ij}$ and $d_{00}$
	terms were assigned independent unresolved signs is replaced here by
	the canonically normalised result of Ref.~\cite{KosteleckyLane1999}.
	
	\subsection{Particle--antiparticle transformation}
	
	A further point is essential for the matter--antimatter analysis.
	For the one-particle SME Hamiltonian, the corresponding antiparticle
	Hamiltonian can be written in the same form after the coefficient
	substitutions
	\begin{equation}
		\begin{aligned}
			a_\mu&\rightarrow -a_\mu,
			&
			b_\mu&\rightarrow +b_\mu,
			&
			c_{\mu\nu}&\rightarrow +c_{\mu\nu},
			\\
			d_{\mu\nu}&\rightarrow -d_{\mu\nu},
			&
			e_\mu&\rightarrow -e_\mu,
			&
			f_\mu&\rightarrow -f_\mu,
			\\
			g_{\lambda\mu\nu}&\rightarrow +g_{\lambda\mu\nu},
			&
			H_{\mu\nu}&\rightarrow -H_{\mu\nu}.
		\end{aligned}
		\label{eq:antiparticleSME}
	\end{equation}
	These substitutions are part of the canonical particle--antiparticle
	construction of the SME Hamiltonian \cite{KosteleckyLane1999}.
	
	Equation~\eqref{eq:antiparticleSME} also illustrates why the statement
	``CPT even'' or ``CPT odd'' should not be used as a shortcut for the
	sign of an experimentally measured energy shift. The discrete-symmetry
	classification concerns the relativistic operator, whereas the
	particle--antiparticle relation of a particular nonrelativistic
	observable follows from the complete Hamiltonian and the conventions
	used for the antiparticle state.
	
	\subsection{Controlled matching summary}
	
	The controlled results are summarised in Table~\ref{tab:matching}.
	The last column deliberately describes operator structure rather than
	assigning a DM potential number to an SME coefficient.
	
	\begin{table*}[t]
		\caption{Controlled nonrelativistic structures obtained from the
			selected minimal-SME coefficients. The table does not identify an SME
			coefficient with a two-body mediator coupling.}
		\label{tab:matching}
		\begin{ruledtabular}
			\begin{tabular}{lcccl}
				SME coefficient & CPT character & Nonrelativistic contribution
				& Leading order & Operator structure\\
				\hline
				$b_i$
				& odd
				& $-b_i\sigma^i$
				& $m^0$
				& spin-linear
				\\
				$b_0$
				& odd
				& $(b_0/m)\boldsymbol{\sigma}\cdot\mathbf p$
				& $m^{-1}$
				& spin--momentum
				\\
				$H_{ij}$
				& even
				& $\mathbf B_H\cdot\boldsymbol{\sigma}$
				& $m^0$
				& spin-linear
				\\
				$H_{0i}$
				& even
				& $-m^{-1}\boldsymbol{\sigma}\cdot
				(\mathbf p\times\mathbf H_E)$
				& $m^{-1}$
				& spin--momentum
				\\
				$d_{i0}$
				& even
				& $m d_{i0}\sigma^i$
				& $m^1$
				& spin-linear
				\\
				$d_{ij},d_{00}$
				& even
				& $-(d_{ij}+d_{00}\delta_{ij})p_j\sigma^i$
				& $m^0$
				& spin--momentum
			\end{tabular}
		\end{ruledtabular}
	\end{table*}
	
	The table is intentionally narrower than a complete SME--DM
	dictionary. The minimal SME fermion sector contains additional
	coefficients, including $a_\mu$, $c_{\mu\nu}$, $e_\mu$, $f_\mu$, and
	$g_{\lambda\mu\nu}$, which require separate treatment. Moreover, a
	true two-body matching requires an explicit relativistic scattering
	amplitude and cannot be inferred from the one-body Hamiltonian alone.
	
	% ============================================================
	\section{Matter--antimatter comparison of constraints}
	% ============================================================
	
	\subsection{Signed observables}
	
	Let $\theta_f$ denote a signed observable associated with a fermion and
	let $\theta_{\bar f}$ be the corresponding observable for the
	antifermion. Once the complete relativistic transformation has been
	derived, exact CPT symmetry imposes a relation of the form
	\begin{equation}
		\theta_{\bar f}
		=
		\eta_\CPT\theta_f,
		\qquad
		\eta_\CPT=\pm1,
		\label{eq:CPTsigned}
	\end{equation}
	where the value of $\eta_\CPT$ depends on the particular signed
	observable and its conventions. It must therefore be derived rather
	than assigned solely from the name of the underlying SME coefficient.
	
	A corresponding CPT residual is
	\begin{equation}
		\Delta_\CPT
		=
		\theta_{\bar f}-\eta_\CPT\theta_f .
		\label{eq:DeltaCPT}
	\end{equation}
	
	If the measurements are approximately Gaussian, a simple test statistic
	is
	\begin{equation}
		\chi^2_\CPT
		=
		\frac{
			\left(
			\theta_{\bar f}-\eta_\CPT\theta_f
			\right)^2
		}{
			\sigma_{\bar f}^2+\sigma_f^2
			-2\eta_\CPT\,\mathrm{Cov}(\theta_f,\theta_{\bar f})
		}.
		\label{eq:chiCPT}
	\end{equation}
	
	This is the appropriate type of construction when signed information
	and its uncertainty are available.
	
	\subsection{One-sided exclusion limits}
	
	The situation is different when the published results are
	\begin{equation}
		|g_f|<U_f,
		\qquad
		|g_{\bar f}|<U_{\bar f},
		\qquad
		U_f,U_{\bar f}>0 .
		\label{eq:limits}
	\end{equation}
	
	One can define the purely numerical quantity
	\begin{equation}
		A_U
		=
		\frac{U_f-U_{\bar f}}
		{U_f+U_{\bar f}},
		\label{eq:AU}
	\end{equation}
	but it should not be called a measured CPT asymmetry.
	
	\subsection{Sensitivity-gap theorem}
	
	Suppose two experiments constrain the same underlying coupling magnitude
	$g_\star$ with sensitivity factors $s_f$ and $s_{\bar f}$:
	\begin{equation}
		U_f=s_f g_\star,
		\qquad
		U_{\bar f}=s_{\bar f}g_\star .
		\label{eq:sensitivity}
	\end{equation}
	
	Substitution into Eq.~\eqref{eq:AU} gives
	\begin{equation}
		A_U
		=
		\frac{s_f-s_{\bar f}}
		{s_f+s_{\bar f}}.
		\label{eq:gapA}
	\end{equation}
	
	The underlying coupling has cancelled. Equation~\eqref{eq:gapA}
	therefore contains no CPT parameter.
	
	For example,
	\begin{equation}
		s_{\bar f}\ll s_f
		\quad\Longrightarrow\quad
		A_U\rightarrow +1,
	\end{equation}
	whereas
	\begin{equation}
		s_f\ll s_{\bar f}
		\quad\Longrightarrow\quad
		A_U\rightarrow -1.
	\end{equation}
	
	The result holds irrespective of whether the underlying signed
	interaction is CPT even or CPT odd. Consequently, a value
	$|A_U|\simeq1$ obtained from two positive exclusion limits cannot by
	itself be interpreted as evidence for CPT violation.
	
	There is also a useful algebraic warning concerning signed
	CPT-odd quantities. If exact CPT symmetry requires
	\begin{equation}
		g_{\bar f}=-g_f,
	\end{equation}
	then the denominator of Eq.~\eqref{eq:Adef} vanishes. Exact
	sign reversal therefore does not correspond to $A=\pm1$. The apparent
	saturation of $A_U$ at $\pm1$ is a property of comparing positive
	magnitudes.
	
	\subsection{Sensitivity-ratio diagnostic}
	
	A more transparent quantity derived from two one-sided limits is
	\begin{equation}
		R_\alpha(\lambda)
		=
		\frac{U_{\bar f,\alpha}(\lambda)}
		{U_{f,\alpha}(\lambda)},
		\label{eq:R}
	\end{equation}
	provided the limits have first been harmonised to the same coupling
	normalisation, confidence convention, and operator definition.
	
	For a symmetric measure of the mismatch one may define
	\begin{equation}
		G_\alpha(\lambda)
		=
		\left|\ln R_\alpha(\lambda)\right|.
		\label{eq:G}
	\end{equation}
	
	Equal sensitivity corresponds to
	\begin{equation}
		R_\alpha=1,
		\qquad
		G_\alpha=0.
	\end{equation}
	
	Unlike $A_U$, $G_\alpha$ makes no claim to measure a signed physical
	coupling. It is explicitly a diagnostic of experimental coverage.
	
	For future experiments that provide signed likelihoods
	$\mathcal L_f(g)$ and $\mathcal L_{\bar f}(g)$, the CPT-conserving
	likelihood can be written schematically as
	\begin{equation}
		\mathcal L_{\CPT}(g)
		=
		\mathcal L_f(g)
		\mathcal L_{\bar f}(\eta_\CPT g),
	\end{equation}
	with the appropriate nuisance parameters and correlations included.
	This is fundamentally different from comparing the amplitudes of two
	exclusion curves.
	
	% ============================================================
	\section{Experimental constraint database}
	% ============================================================
	
	\subsection{Compilation strategy}
	
	To quantify the present experimental landscape, we constructed a
	structured registry of published constraints on exotic spin-dependent
	interactions. Each record contains, where available, the fermion
	sectors, coupling family, DM potential label, interaction range, and
	published constraint. The common interaction-range coordinate is
	\begin{equation}
		\lambda=\frac{\hbar}{m_\phi c}.
	\end{equation}
	
	The registry contains
	\begin{equation}
		N_{\rm compiled}=273,
		\qquad
		N_{\rm analysed}=258.
		\label{eq:Nrecords}
	\end{equation}
	
	Fifteen records are held out of the physics analysis. Twelve are
	combined $g_Ag_V$ envelope curves labelled in the source database as
	combined or review-only and lacking sufficient information for a unique
	potential assignment. Two are contributor-template placeholders. One
	is an astrophysical electron--nucleon bound for which the available
	source information does not establish a unique potential assignment.
	
	The analysed records are distributed among the principal coupling
	families as shown in Table~\ref{tab:families}.
	
	\begin{table}[t]
		\caption{Constraint records in the compiled registry. The counts are
			registry entries, not statistically independent measurements.}
		\label{tab:families}
		\begin{ruledtabular}
			\begin{tabular}{lrr}
				Coupling family & Compiled & Analysed\\
				\hline
				$g_Ag_A$ & 63 & 63\\
				$g_Pg_S$ & 53 & 53\\
				$g_Sg_S$ & 49 & 49\\
				$g_Vg_V$ & 43 & 43\\
				$g_Ag_V$ & 37 & 25\\
				$g_Pg_P$ & 15 & 14\\
				$V_1$ unlabelled coupling & 11 & 11\\
				Template placeholder & 2 & 0\\
				\hline
				Total & 273 & 258
			\end{tabular}
		\end{ruledtabular}
	\end{table}
	
	The registry includes constraints from atomic and molecular
	spectroscopy, torsion balances, comagnetometry, spin sensors,
	positronium, antiprotonic helium, and astrophysical observations.
	These platforms probe substantially different interaction ranges and
	spin sectors \cite{Ficek2018,Adkins2022,Liang2023,Jiang2024,
		Cong2025,WuYan2026}.
	
	\subsection{Data-quality controls}
	
	The numerical registry requires several qualifications.
	
	First, records may be filed under more than one coupling family.
	Consequently, the 258 analysed records do not represent 258
	statistically independent measurements. Forty-eight records correspond
	to repeated entries of the same bound under different coupling labels,
	leaving 210 unique constraint curves for the purposes of the curve
	inventory. The term ``unique'' is used deliberately: these curves
	should not be treated as statistically independent unless their
	experimental origins are independently established.
	
	Second, six registry entries contain interaction-range values of order
	$10^{52}\,\mathrm{m}$. These values are physically implausible and are
	consistent with a unit-token parsing error. They are excluded from
	interpretations of global range coverage and do not enter the direct
	matter--antimatter comparisons discussed below.
	
	Third, several experiments report combined potential classes rather
	than individual components. Labels such as $V_{2+3}$ and $V_{4+5}$ are
	therefore retained as published. They are not decomposed into separate
	potential limits without an independent theoretical justification.
	
	These controls are important because a compilation of exclusion curves
	can otherwise create an artificial impression of statistical precision.
	
	% ============================================================
	\section{Matter--antimatter coverage}
	% ============================================================
	
	Of the 258 analysed records, 19 involve an antimatter sector:
	nine electron--antiproton records, five electron--positron records,
	three electron--antimuon records, one antiproton--helium record, and
	one molecular-ion record associated with an antimuon sector. Thus
	\begin{equation}
		N_{\rm antimatter}=19,
		\qquad
		N_{\rm matter}=239.
		\label{eq:mattercounts}
	\end{equation}
	
	The relevant issue for a CPT comparison is not simply whether an
	antimatter measurement exists. A directly comparable result requires
	agreement in operator definition, fermion sector, coupling
	normalisation, and overlapping interaction range.
	
	The sector inventory in the present registry is shown in
	Table~\ref{tab:sectorcoverage}. The sector tags are not necessarily
	additive because a single registry record may contribute to more than
	one phenomenological category.
	
	\begin{table*}[t]
		\caption{Matter and antimatter coverage by the canonical fermion-sector
			labels used in the registry. Counts are analysed registry records and
			are not statistically independent measurements. Sector tags are
			non-exclusive.}
		\label{tab:sectorcoverage}
		\begin{ruledtabular}
			\begin{tabular}{lrrrr}
				Matter sector & Antimatter sector & Matter & Antimatter
				& Coverage\\
				\hline
				$e-e$ & $e-e^+$ & 44 & 5 & both\\
				$e-p$ & $e-\bar p$ & 12 & 9 & both\\
				$e-\mu$ & $e-\bar\mu$ & 3 & 3 & both; not directly matched\\
				$e-n$ & $e-\bar n$ & 16 & 0 & no antimatter record\\
				$\mu-\mu$ & $\mu-\bar\mu$ & 1 & 0 & no antimatter record\\
				$n-p$ & $n-\bar p$ & 9 & 0 & no directly matched record\\
				$n-n$ & $n-\bar n$ & 8 & 0 & no antimatter record\\
				$p-p$ & $p-\bar p$ & 0 & 0 & no record\\
				$e-N$ & $e-\bar N$ & 54 & 0 & no directly matched record\\
				$n-N$ & $n-\bar N$ & 52 & 0 & no directly matched record\\
				$p-N$ & $p-\bar N$ & 23 & 0 & no directly matched record
			\end{tabular}
		\end{ruledtabular}
	\end{table*}
	
	The proton--antinucleus entry requires particular care. Antiprotonic
	helium spectroscopy provides an antiproton--nuclear constraint
	\cite{Ficek2018}, but the published result uses a normalisation that
	cannot be inserted into the corresponding proton--nucleus entries
	without an explicit conversion. It is therefore not counted as a
	directly comparable $p-\bar N$ result in the adopted coupling
	normalisation.
	
	The three nuclear-source categories contain
	\begin{equation}
		54+52+23=129
	\end{equation}
	matter-sector registry records. The corresponding directly comparable
	antimatter count is zero in the adopted convention. The principal
	coverage limitation is therefore not the absence of matter-sector
	constraints but the lack of matched antimatter measurements in sectors
	where matter data are already numerous.
	
	% ============================================================
	\section{Direct matter--antimatter comparisons}
	% ============================================================
	
	\subsection{Matched pairs}
	
	Applying the registry matching criteria to coupling family, potential
	label, fermion sector, and overlapping interaction range gives eleven
	direct matter--antimatter pairs.
	
	For each pair, the published curves were interpolated onto a common
	logarithmic range coordinate. A nominal 300-point grid was used in the
	original numerical implementation. The pointwise quantity was
	\begin{equation}
		A_U(\lambda_i)
		=
		\frac{
			U_f(\lambda_i)-U_{\bar f}(\lambda_i)
		}{
			U_f(\lambda_i)+U_{\bar f}(\lambda_i)
		}.
		\label{eq:AUlambda}
	\end{equation}
	
	The values in Table~\ref{tab:asymmetry} are retained as diagnostic
	quantities only. They are not interpreted as estimates of a CPT
	parameter.
	
	\begin{table*}[t]
		\caption{Bound-based matter--antimatter sensitivity diagnostics for the
			eleven directly matched pairs in the registry. The quantity $|A_U|$
			is calculated from positive upper limits and is therefore not a
			CPT-violation estimator. The effective degrees of freedom are included
			only to document the exploratory autocorrelation analysis. The
			previously quoted bootstrap intervals are not reproduced because the
			underlying uncertainty model does not provide formal confidence
			intervals for a CPT parameter.}
		\label{tab:asymmetry}
		\begin{ruledtabular}
			\begin{tabular}{lccc}
				Coupling & Potential & Sector & $|A_U|$ \\
				\hline
				$g_Ag_A$ & $V_2$ & $e-\bar p$ & 0.9998\\
				$g_Ag_A$ & $V_{2+3}$ & $e-e^+$ & 0.9892\\
				$g_Ag_A$ & $V_{2+3}$ & $e-e^+$ & 0.9539\\
				$g_Pg_P$ & $V_{2+3}$ & $e-e^+$ & 0.9535\\
				$g_Vg_V$ & $V_{2+3}$ & $e-e^+$ & 0.9535\\
				$g_Sg_S$ & $V_1$ & $e-e^+$ & 0.8727\\
				$g_Ag_A$ & $V_{2+3}$ & $e-\bar p$ & 0.8236\\
				$g_Ag_A$ & $V_{2+3}$ & $e-\bar p$ & 0.8044\\
				$g_Pg_P$ & $V_{2+3}$ & $e-\bar p$ & 0.7994\\
				$g_Vg_V$ & $V_{2+3}$ & $e-\bar p$ & 0.7994\\
				$g_Ag_A$ & $V_2$ & $e-e^+$ & 0.3336
			\end{tabular}
		\end{ruledtabular}
	\end{table*}
	
	The eleven values range from approximately $0.33$ to $1$. This broad
	range is precisely what is expected if the statistic is responding to
	different experimental sensitivities rather than to a common signed
	physical asymmetry.
	
	\subsection{Internal null behaviour}
	
	The $g_Sg_S$, $V_1$, electron--positron comparison is particularly
	instructive because it produces a large bound-based asymmetry,
	$|A_U|\simeq0.87$, despite the fact that the quantity being compared is
	not a signed CPT-odd observable. The same procedure therefore produces
	large apparent asymmetries in channels where interpreting the result as
	a CPT signal is unjustified.
	
	This is an internal null check on the statistic. The numerical value of
	$A_U$ is controlled by the relative strengths of the two exclusion
	curves, not by the discrete-symmetry classification of the underlying
	interaction.
	
	The lowest value in the table,
	\begin{equation}
		|A_U|\simeq0.334,
	\end{equation}
	also demonstrates that the statistic is not mathematically forced to
	be close to unity. Its magnitude simply tracks the degree of mismatch
	between the available constraints.
	
	% ============================================================
	\section{Coverage by potential label}
	% ============================================================
	
	The same registry can be organised by its DM potential labels. The
	result is shown in Table~\ref{tab:potentialcoverage}. The counts in
	different rows are not additive because individual records can carry
	multiple phenomenological labels.
	
	\begin{table}[t]
		\caption{Matter and antimatter registry coverage by DM potential label.
			Counts are non-exclusive registry counts.}
		\label{tab:potentialcoverage}
		\begin{ruledtabular}
			\begin{tabular}{lrr}
				Potential label & Matter & Antimatter\\
				\hline
				$V_{4+5}$ & 58 & 3\\
				$V_{9+10}$ & 31 & 0\\
				$V_{2+3}$ & 27 & 10\\
				$V_{1a}$ & 30 & 0\\
				$V_2$ & 17 & 2\\
				$V_1$ & 44 & 3\\
				$V_{12+13}$ & 16 & 0\\
				$V_{11}$ & 7 & 0\\
				$V_8$ & 5 & 1\\
				$V_{15}$ & 2 & 0\\
				$V_{16}$ & 2 & 0
			\end{tabular}
		\end{ruledtabular}
	\end{table}
	
	Six of the eleven classified potential groups have no antimatter record
	despite nonzero matter coverage:
	\begin{equation}
		V_{11},\qquad
		V_{15},\qquad
		V_{16},\qquad
		V_{1a},\qquad
		V_{9+10},\qquad
		V_{12+13}.
	\end{equation}
	
	The $V_{2+3}$ group has the largest number of antimatter entries in the
	present registry. This is a statement about the existing experimental
	literature and not about an intrinsic suitability of this potential for
	a CPT test.
	
	Recent experimental developments continue to improve sensitivity to
	spin-dependent interactions using quantum sensors and related
	techniques \cite{Jiao2021,Liang2023,Jiang2024,Cong2025,WuYan2026}.
	Recent torsion-balance work has also explicitly considered
	$V_{4+5}$ and $V_{12+13}$ interactions between polarized electrons and
	unpolarised nucleons \cite{Jin2026}. These developments demonstrate
	that improved matter-sector coverage is continuing, while the present
	registry still shows a substantial antimatter deficit.
	
	% ============================================================
	\section{Statistical methodology and limitations}
	% ============================================================
	
	\subsection{Interpolation and correlation}
	
	The 300-point interpolation grid used for the matched curves should not
	be interpreted as 300 independent observations. Published exclusion
	curves are smooth functions of interaction range and neighbouring
	interpolated points are strongly correlated.
	
	The exploratory analysis estimated the autocorrelation function of the
	interpolated residual sequence,
	\begin{equation}
		C(k)
		=
		\frac{
			\sum_i(x_i-\bar x)(x_{i+k}-\bar x)
		}{
			\sum_i(x_i-\bar x)^2
		},
		\label{eq:ACF}
	\end{equation}
	and used the first crossing of $C(k)=e^{-1}$ as a characteristic
	correlation scale.
	
	This procedure is useful as a diagnostic against treating every
	interpolated point as independent. It is not a rigorous derivation of
	the likelihood of the original experimental measurements, because the
	published curves do not generally provide the full covariance matrix.
	
	Accordingly, the autocorrelation-based effective degrees of freedom are
	not used as the primary basis for a CPT significance claim.
	
	\subsection{Uncertainty assignment}
	
	Most published constraint curves do not provide point-by-point
	covariance matrices. The original analysis therefore estimated local
	fractional uncertainties from the curvature of the curves in
	logarithmic coordinates. A schematic weighted diagnostic is
	\begin{equation}
		\chi^2_{\rm w}
		=
		\sum_i
		\frac{
			[A_U(\lambda_i)-\overline{A_U}]^2
		}{
			\sigma_A^2(\lambda_i)
		}.
		\label{eq:chiweighted}
	\end{equation}
	
	A parametric bootstrap with 2000 resamples was then used to propagate
	the estimated uncertainty model to summary quantities. This is not a
	bootstrap of the original experimental events. It propagates an
	assumed uncertainty model inferred from published curves and therefore
	cannot restore systematic uncertainties or correlations that were not
	published.
	
	Furthermore, taking an absolute value produces a nonlinear statistic.
	An interval for $|A_U|$ therefore cannot be interpreted as a confidence
	interval for a signed CPT parameter.
	
	\subsection{The correct null hypothesis}
	
	A conventional signed CPT test asks whether
	\begin{equation}
		H_0:
		\theta_{\bar f}=\eta_\CPT\theta_f.
		\label{eq:nullCPT}
	\end{equation}
	
	The one-sided-bound comparison instead asks whether
	\begin{equation}
		H_0:
		\frac{U_{\bar f}}{U_f}=1.
		\label{eq:nullsens}
	\end{equation}
	
	These hypotheses are not equivalent. Consequently, even a very precise
	determination that $U_{\bar f}/U_f\neq1$ is not a measurement of
	CPT violation.
	
	% ============================================================
	\section{A practical framework for future CPT tests}
	% ============================================================
	
	A genuine matter--antimatter test of an exotic spin-dependent
	interaction requires four conditions.
	
	First, the same physical operator must be defined on both sides of the
	comparison. Spin, momentum, velocity, and antiparticle-state
	conventions must be specified explicitly.
	
	Second, the transformation between the matter and antimatter
	observables must be derived from the relativistic operator or
	scattering amplitude. The sign of an effective nonrelativistic
	Hamiltonian term cannot safely be inferred from a coupling label alone.
	
	Third, the experiments must provide signed information or likelihood
	functions whenever a signed CPT test is intended. Two positive upper
	limits do not determine whether the underlying coupling is positive,
	negative, or zero.
	
	Fourth, the matter and antimatter experiments must have sufficiently
	overlapping sensitivity in interaction range and operator convention.
	A comparison made where one experiment is many orders of magnitude
	less sensitive than the other is principally a comparison of
	experimental reach.
	
	\subsection{Stage I: sensitivity coverage}
	
	For each operator and range, calculate
	\begin{equation}
		R_\alpha(\lambda)
		=
		\frac{U_{\bar f,\alpha}(\lambda)}
		{U_{f,\alpha}(\lambda)}.
	\end{equation}
	
	A balanced interval may be defined by
	\begin{equation}
		|\ln R_\alpha(\lambda)|<\epsilon,
		\label{eq:balanced}
	\end{equation}
	where $\epsilon$ is selected according to the intended experimental
	criterion.
	
	This step identifies where matter and antimatter measurements have
	comparable sensitivity. It does not itself test CPT.
	
	\subsection{Stage II: likelihood-level comparison}
	
	Suppose a matter experiment provides a signed estimator $\hat g_f$ and
	likelihood $\mathcal L_f(g_f)$, while the corresponding antimatter
	experiment provides $\mathcal L_{\bar f}(g_{\bar f})$. The
	CPT-conserving hypothesis is
	\begin{equation}
		g_{\bar f}=\eta_\CPT g_f.
	\end{equation}
	
	A likelihood-ratio statistic can then be defined as
	\begin{equation}
		q_\CPT
		=
		-2\ln
		\frac{
			\displaystyle
			\max_g
			\mathcal L_f(g)
			\mathcal L_{\bar f}(\eta_\CPT g)
		}{
			\displaystyle
			\max_{g_f,g_{\bar f}}
			\mathcal L_f(g_f)
			\mathcal L_{\bar f}(g_{\bar f})
		}.
		\label{eq:likelihoodratio}
	\end{equation}
	
	In an actual analysis, nuisance parameters, correlated systematics,
	non-Gaussian likelihoods, and the precise statistical construction of
	each experiment must be retained. Equation~\eqref{eq:likelihoodratio}
	is therefore a framework rather than a numerical result of the present
	database.
	
	% ============================================================
	\section{Experimental implications}
	% ============================================================
	
	\subsection{Electron--nucleus, neutron--nucleus, and proton--nucleus
		sectors}
	
	The registry contains substantial matter-sector coverage in
	\begin{equation}
		N_{e-N}=54,
		\qquad
		N_{n-N}=52,
		\qquad
		N_{p-N}=23,
	\end{equation}
	while the corresponding directly comparable antimatter counts are
	zero in the adopted coupling normalisation.
	
	These sectors consequently provide clear targets for future matched
	matter--antimatter measurements. The objective should not simply be a
	stronger upper limit. The experiment should reproduce the same
	operator definition, coupling normalisation, and interaction-range
	dependence as an existing matter-sector measurement so that the two
	results can eventually be compared at the likelihood level.
	
	\subsection{Velocity-dependent and related potential classes}
	
	The registry contains matter-sector entries but no antimatter entries
	for $V_{1a}$, $V_{9+10}$, and $V_{12+13}$. Such gaps identify areas
	where a future antimatter measurement could add qualitatively new
	coverage.
	
	The recent torsion-balance proposal of Jin \textit{et al.} specifically
	addresses $V_{4+5}$ and $V_{12+13}$ using polarized electrons and
	unpolarised nucleon sources \cite{Jin2026}. The work illustrates how
	operator-specific sensitivity can be developed experimentally, but a
	corresponding antimatter source or bound system would still be required
	for a direct CPT comparison.
	
	\subsection{Electron--positron and electron--antiproton systems}
	
	Electron--positron and electron--antiproton systems currently provide
	the most developed antimatter coverage in the registry. Positronium is
	particularly useful because its leptonic composition allows precision
	comparison with bound-state QED calculations \cite{Adkins2022}.
	Atomic spectroscopy has also been used to constrain exotic
	spin-dependent interactions in positronium, muonium, hydrogen, and
	antiprotonic helium \cite{Fadeev2022}.
	
	Antiprotonic helium provides a direct example of a matter--antimatter
	spin-dependent constraint, having been used to search for exotic
	spin- and velocity-dependent interactions between electrons and
	antiprotons \cite{Ficek2018}.
	
	These systems are therefore useful benchmark channels for developing
	the signed-likelihood methodology advocated here. Where physically
	possible, future measurements should preserve enough information to
	construct signed shifts, likelihood functions, or covariance matrices
	rather than reducing the final result to a positive exclusion magnitude.
	
	% ============================================================
	\section{Discussion}
	% ============================================================
	
	The main result of this work is methodological. A numerical difference
	between matter and antimatter exclusion limits is not, by itself, a
	measurement of CPT violation.
	
	The result follows directly from Eq.~\eqref{eq:gapA}. If two experiments
	report positive upper limits, their difference can arise entirely from
	relative experimental sensitivity. This remains true even if the
	underlying interaction is exactly CPT conserving.
	
	The compiled registry gives this mathematical observation an empirical
	context. Only 19 of 258 analysed records involve an antimatter sector,
	and six of eleven classified potential groups have no antimatter
	coverage. The largest sector deficits occur in electron--nucleus,
	neutron--nucleus, and proton--nucleus interactions.
	
	The numerical values of $|A_U|$ should consequently be described as
	sensitivity-gap diagnostics. Values in the range
	\begin{equation}
		|A_U|\simeq0.8--1
	\end{equation}
	do not constitute anomalous CPT signals merely because they are
	numerically large.
	
	The reinterpretation does not diminish the value of antimatter
	experiments. Instead, it specifies what additional information is
	required before their results can become quantitative CPT tests.
	The experimentally useful objective is to produce matched matter and
	antimatter measurements of the same operator over overlapping
	interaction ranges.
	
	The recent literature demonstrates continued progress on both the
	matter and antimatter sides. Spin sensors and NV-diamond experiments
	have improved laboratory constraints \cite{Jiao2021,Liang2023,Jiang2024},
	while antiprotonic helium and positronium provide antimatter systems
	with well-developed precision spectroscopy
	\cite{Ficek2018,Adkins2022}. New torsion-balance proposals extend the
	available sensitivity to additional spin-dependent structures
	\cite{Jin2026}. The principal missing ingredient for a broad CPT
	programme is therefore not another numerical comparison of unequal
	upper limits, but matched measurements with compatible definitions.
	
	There is also a broader methodological lesson for global compilations
	of new-physics constraints. A collection of exclusion curves is not
	automatically a collection of statistically independent measurements.
	Duplicate records, combined potential classes, correlated interpolation
	points, different coupling normalisations, and heterogeneous
	statistical conventions must be retained explicitly.
	
	% ============================================================
	\section{Limitations}
	% ============================================================
	
	The numerical registry is a curated compilation rather than a guarantee
	of exhaustive coverage of every relevant publication available by the
	date of this manuscript. The literature continues to develop rapidly,
	and future releases should incorporate additional measurements and
	proposals.
	
	Fifteen of the 273 compiled records are held out of the physics
	analysis. This is preferable to assigning them unsupported potential
	labels, but it means that the registry is not completely closed under
	all possible source conventions.
	
	Several range entries are affected by probable unit-token errors. These
	entries must be corrected against their source data before the registry
	is used for global range statistics.
	
	The SME--nonrelativistic matching is restricted to
	$b_\mu$, $H_{\mu\nu}$, and $d_{\mu\nu}$. It is therefore not a complete
	SME fermion-sector dictionary. In particular, the remaining minimal-SME
	coefficients require separate derivations.
	
	The DM labels in the database are phenomenological labels and are not
	treated as one-to-one images of SME coefficients. A complete
	relativistic--nonrelativistic--two-body dictionary would require explicit
	matching of relativistic scattering amplitudes for the underlying
	mediator models.
	
	Finally, the bound-asymmetry analysis is not a CPT hypothesis test.
	The autocorrelation analysis can diagnose the effect of correlated
	interpolation, but it cannot convert one-sided exclusion curves into
	signed measurements or reconstruct experimental likelihoods that were
	not published.
	
	% ============================================================
	\section{Conclusions}
	% ============================================================
	
	We have developed a unified framework for connecting selected
	relativistic SME fermion operators to their controlled
	nonrelativistic structures and for assessing what can and cannot be
	learned from matter--antimatter comparisons of published
	spin-dependent interaction limits.
	
	The canonical FW treatment gives, through the leading orders relevant
	here,
	\begin{align}
		H_{\NR}^{(b)}
		&=
		-b_i\sigma^i
		+\frac{b_0}{m}\,
		\boldsymbol{\sigma}\cdot\mathbf p,
		\label{eq:summaryb}
		\\
		H_{\NR}^{(H)}
		&=
		\mathbf B_H\cdot\boldsymbol{\sigma}
		-\frac{1}{m}
		\boldsymbol{\sigma}\cdot
		(\mathbf p\times\mathbf H_E),
		\label{eq:summaryH}
		\\
		H_{\NR}^{(d)}
		&=
		m d_{i0}\sigma^i
		-
		(d_{ij}+d_{00}\delta_{ij})p_j\sigma^i .
		\label{eq:summaryd}
	\end{align}
	
	The $d_{i0}$ contribution is mass enhanced and agrees with the
	established nonrelativistic SME Hamiltonian. The momentum-linear
	$d$ sector is governed by the canonically normalised combination
	$d_{ij}+d_{00}\delta_{ij}$; the earlier treatment in terms of separate
	unresolved signs is therefore superseded.
	
	The particle--antiparticle relation must likewise be obtained from the
	complete Hamiltonian. In the canonical SME treatment the relevant
	coefficient substitutions include
	\begin{equation}
		b_\mu\rightarrow+b_\mu,
		\qquad
		d_{\mu\nu}\rightarrow-d_{\mu\nu},
		\qquad
		H_{\mu\nu}\rightarrow-H_{\mu\nu},
	\end{equation}
	for the antiparticle Hamiltonian \cite{KosteleckyLane1999}. The CPT
	classification of a relativistic coefficient should not be confused
	with the sign of an individual experimentally defined nonrelativistic
	observable.
	
	The principal statistical result is exact. If two positive upper
	limits are represented by
	\begin{equation}
		U_f=s_f g_\star,
		\qquad
		U_{\bar f}=s_{\bar f}g_\star,
	\end{equation}
	then
	\begin{equation}
		A_U
		=
		\frac{s_f-s_{\bar f}}
		{s_f+s_{\bar f}}.
	\end{equation}
	The underlying coupling cancels, so the statistic is a measure of
	relative experimental sensitivity rather than a CPT observable.
	Large values of $|A_U|$ can occur under exact CPT symmetry, while
	small values do not by themselves establish CPT invariance.
	
	The compiled registry contains 273 records, 258 of which enter the
	present analysis. Nineteen involve antimatter, and eleven satisfy the
	adopted criteria for direct matter--antimatter comparison. Six of the
	eleven classified potential groups have no antimatter record. The
	largest sector deficits occur in electron--nucleus, neutron--nucleus,
	and proton--nucleus interactions.
	
	The appropriate next-generation programme is therefore two-stage.
	First, existing exclusion curves should be used to identify balanced
	sensitivity regions through
	\begin{equation}
		R_\alpha(\lambda)
		=
		\frac{U_{\bar f,\alpha}(\lambda)}
		{U_{f,\alpha}(\lambda)}.
	\end{equation}
	Second, experiments in those regions should report signed observables
	or likelihood functions so that the actual CPT hypothesis
	\begin{equation}
		\theta_{\bar f}-\eta_\CPT\theta_f=0
	\end{equation}
	can be tested directly.
	
	% ============================================================
	\section{Further Work}
	% ============================================================
	
	Several extensions follow naturally.
	
	First, the complete minimal-SME fermion sector should be treated using
	the same canonical field-normalisation and FW procedure. This would
	extend the present controlled matching beyond $b_\mu$, $H_{\mu\nu}$,
	and $d_{\mu\nu}$.
	
	Second, the one-body SME structures should be matched explicitly to
	two-fermion relativistic amplitudes for specified mediator models.
	This is the appropriate route to a genuine SME--DM coupling dictionary.
	
	Third, the constraint registry should preserve source likelihoods,
	confidence conventions, covariance information, coupling
	normalisations, and experimental provenance whenever these are
	available. Such a structure would allow future matter and antimatter
	measurements to be compared without reconstructing the statistical
	problem from published exclusion curves.
	
	Finally, the most informative future experiments are those that produce
	matched matter--antimatter observables for the same operator over
	overlapping interaction ranges. Once such measurements exist, the
	sensitivity-ratio analysis developed here can identify the relevant
	coverage, while a likelihood-level comparison can provide the actual
	quantitative CPT test.
	
	% ============================================================
	\section*{Acknowledgments}
	% ============================================================
	
	We thank the experimental and theoretical communities whose published
	results and constraint data provide the foundation for systematic
	cross-platform comparisons of exotic spin-dependent interactions.
	TSO gratefully acknowledges CERN for sponsoring his participation in
	the CERN Summer Student Lecture Programme 2026 and for the scientific
	training and research environment provided during his participation.
	He also acknowledges the CMS Collaboration for the broader scientific
	training and collaborative research environment which supported the development of this work.
	
	% ============================================================
	% References
	% ============================================================
	
	\bibliographystyle{apsrev4-2}
	\bibliography{references}
	
\end{document}